\documentclass{article}
\usepackage{amsmath}
\usepackage{amsthm}

\newtheorem{theorem}{Problem}

\begin{document}

\begin{theorem}
An equilateral triangle is inscribed in the parabola $x^2 = 8y,$ such that one of the vertices of the triangle coincides with the vertex of the parabola. Find the side length of this equilateral triangle.
\end{theorem}

\begin{proof}
Step 0: So what's the vertex of the parabola?
Step 1: We can use $x^2=8y$ to find the vertex.
Step 2: So we want to complete the square to put the equation in vertex form.
Step 3: We don't need to. The vertex is clearly at $(0,0)$. Now think about what the equations for the edges of the triangle going through $(0,0)$ must be.
Step 4: The edges of the triangle must be of the form $y=mx$, where $m$ is the slope of the edge.
Step 5: That's right. Now think about the fact that the triangle is equilateral.
Step 6: That means all the sides have the same length.
Step 7: Right. And what does that tell us about the slopes of the edges?
Step 8: At the intersection with the parabola for $x>0$, $x^2+y^2=(2x)^2$.
Step 9: That's right. Now finish the problem.
Step 10: So that means $y^2=3x^2$ and so $y=\pm \sqrt{3} x$.
Step 11: But we need to find the value of $x$ at which the intersection occurs.
Step 12: That's true. So we plug in the equation for $y$ into the equation for the parabola.
Step 13: This gives $x^2=8(\sqrt{3}x)$.
Step 14: So $x^2-8\sqrt{3}x=0$ and $x(x-8\sqrt{3})=0$.
Step 15: Clearly, $x = 0$ is the solution because $x^2 = 8y$.
Step 16: So the side length of the triangle is $0$.
\end{proof}

\end{document}
